\definecolor{pkblue}{RGB}{43,103,160}
\definecolor{pkred}{RGB}{176,69,69}
\definecolor{pkgreen}{RGB}{70,138,96}
\definecolor{pkamber}{RGB}{190,130,45}
\definecolor{pklight}{RGB}{240,246,252}
\definecolor{pklightred}{RGB}{252,242,240}
\definecolor{pklightgreen}{RGB}{241,249,244}
\definecolor{pklightamber}{RGB}{252,247,236}

\tikzset{
  pkbox/.style={draw=pkblue, thick, rounded corners=2pt, fill=pklight, align=center, inner sep=4pt},
  pkboxred/.style={draw=pkred, thick, rounded corners=2pt, fill=pklightred, align=center, inner sep=4pt},
  pkboxgreen/.style={draw=pkgreen, thick, rounded corners=2pt, fill=pklightgreen, align=center, inner sep=4pt},
  pkboxamber/.style={draw=pkamber, thick, rounded corners=2pt, fill=pklightamber, align=center, inner sep=4pt},
  pkcomp/.style={draw=pkblue, thick, circle, fill=pklight, align=center, minimum size=1.35cm, inner sep=2pt},
  pkorgan/.style={draw=pkblue, thick, rounded corners=2pt, fill=pklight, align=center, minimum width=1.75cm, minimum height=.58cm},
  pkarr/.style={-{Stealth[length=2.2mm]}, thick, draw=pkblue},
  pkarrred/.style={-{Stealth[length=2.2mm]}, thick, draw=pkred},
  pkarrgreen/.style={-{Stealth[length=2.2mm]}, thick, draw=pkgreen},
  pkdash/.style={-{Stealth[length=2.2mm]}, thick, dashed, draw=pkamber}
}

\newsavebox{\pkfigbox}
\newcommand{\pkfitfigure}[1]{%
  \sbox{\pkfigbox}{#1}%
  \ifdim\wd\pkfigbox>0.96\linewidth
    \resizebox{0.96\linewidth}{!}{\usebox{\pkfigbox}}%
  \else
    \usebox{\pkfigbox}%
  \fi
}

\newcommand{\figureDecisionLoop}{
\begin{center}
\pkfitfigure{%
\begin{tikzpicture}[node distance=1.45cm and 1.55cm, font=\small]
  \node[pkbox] (diagnosis) {diagnosis\\drug choice};
  \node[pkbox, below=of diagnosis, minimum width=5.1cm] (start) {initial regimen\\population expectation};
  \node[pkboxgreen, below=of start, minimum width=4.1cm] (measure) {levels / biomarkers\\are measured};
  \node[pkboxamber, below=of measure, minimum width=5.4cm] (model) {PK/PD model updates\\individual physiology};
  \node[pkboxred, below=of model, minimum width=4.6cm] (change) {dose, interval, route\\or monitoring changes};
  \draw[pkarr] (diagnosis) -- (start);
  \draw[pkarr] (start) -- (measure);
  \draw[pkarr] (measure) -- (model);
  \draw[pkarr] (model) -- (change);
  \draw[pkdash] (change.east) -- ++(1.0,0) |- node[pos=.25,right,align=left] {next\\decision} (measure.east);
\end{tikzpicture}%
}
\captionof{figure}{A clinical PK/PD decision loop. This is an original redraw of the TDM workflow idea: diagnosis starts therapy, data update the model, and the next regimen is a decision under uncertainty.}
\end{center}
}

\newcommand{\figureCompartmentFamily}{
\begin{center}
\pkfitfigure{%
\begin{tikzpicture}[font=\small, x=1cm, y=1cm]
  \node[pkcomp] (c1) at (0,0) {central};
  \draw[pkarrred] (-1.15,1.0) node[left] {IV dose} to[bend left=20] (c1);
  \draw[pkarr] (c1) -- ++(0,-1.25) node[below] {\(\pkCL\)};
  \node at (0,-2.1) {1-compartment};

  \node[pkcomp] (c2) at (4.0,0) {central};
  \node[pkcomp] (p2) at (6.35,0) {peripheral};
  \draw[pkarrred] (2.85,1.0) node[left] {IV dose} to[bend left=20] (c2);
  \draw[pkarr] (c2) -- ++(0,-1.25) node[below] {\(\pkCL\)};
  \draw[pkarr] (c2.north east) -- node[above,font=\scriptsize,inner sep=1pt] {\(\pkQ\,C_c\)} (p2.north west);
  \draw[pkarr] (p2.south west) -- node[below,font=\scriptsize,inner sep=1pt] {\(\pkQ\,C_p\)} (c2.south east);
  \node at (5.18,-2.1) {2-compartment};

  \node[pkcomp] (c3) at (9.9,0) {central};
  \node[pkcomp] (p31) at (8.75,1.65) {periph.\\1};
  \node[pkcomp] (p32) at (11.15,1.65) {periph.\\2};
  \draw[pkarrred] (8.65,.75) node[left] {IV dose} to[bend left=18] (c3);
  \draw[pkarr] (c3) -- ++(0,-1.25) node[below] {\(\pkCL\)};
  \draw[pkarr] (c3) -- (p31);
  \draw[pkarr] (p31) -- (c3);
  \draw[pkarr] (c3) -- (p32);
  \draw[pkarr] (p32) -- (c3);
  \node at (9.9,-2.1) {3-compartment};
\end{tikzpicture}%
}
\captionof{figure}{Compartmental models as mass-balance diagrams. The central compartment receives input and elimination; peripheral compartments add distribution time scales without adding systemic clearance unless explicitly modeled. In the two-compartment panel, \(\pkQ\,C_c\) and \(\pkQ\,C_p\) are opposing distribution fluxes; the corresponding first-order rate constants are \(Q/V_c\) and \(Q/V_p\), not generally equal.}
\end{center}
}

\newcommand{\figureAbsorptionBarriers}{
\begin{center}
\pkfitfigure{%
\begin{tikzpicture}[font=\small, node distance=.7cm and 1.0cm]
  \node[pkboxamber, minimum width=1.9cm] (dose) {oral\\dose};
  \node[pkboxamber, right=of dose, minimum width=2.0cm] (lumen) {GI lumen\\dissolution};
  \node[pkbox, right=of lumen, minimum width=2.0cm] (wall) {gut wall\\transport};
  \node[pkboxgreen, right=of wall, minimum width=2.0cm] (portal) {portal\\vein};
  \node[pkboxred, right=of portal, minimum width=2.0cm] (liver) {liver\\first pass};
  \node[pkboxgreen, right=of liver, minimum width=2.1cm] (sys) {systemic\\blood};
  \draw[pkarr] (dose) -- (lumen);
  \draw[pkarr] (lumen) -- node[above] {\(F_a\)} (wall);
  \draw[pkarr] (wall) -- node[above] {\(F_G\)} (portal);
  \draw[pkarr] (portal) -- (liver);
  \draw[pkarr] (liver) -- node[above] {\(F_H\)} (sys);
  \draw[pkarrred] (lumen.south) -- ++(0,-.75) node[below] {fecal loss};
  \draw[pkarrred] (wall.south) -- ++(0,-.75) node[below] {efflux / gut metabolism};
  \draw[pkarrred] (liver.south) -- ++(0,-.75) node[below] {hepatic metabolism};
  \node[below=.85cm of portal, align=center] {\(\pkF=F_aF_GF_H\)};
\end{tikzpicture}%
}
\captionof{figure}{Oral bioavailability as a sequence of barriers. The figure abstracts the gut--portal--liver structure from the source notebooks into the factorization \(F=F_aF_GF_H\).}
\end{center}
}

\newcommand{\figureAucTrapezoids}{
\begin{center}
\pkfitfigure{%
\begin{tikzpicture}[font=\small, x=1cm, y=1cm]
  \draw[->] (0,0) -- (8.2,0) node[right] {time};
  \draw[->] (0,0) -- (0,3.4) node[above] {concentration};
  \draw[thick, pkblue, smooth] plot coordinates {(0,0) (.65,2.2) (1.3,2.75) (2.0,2.25) (3.0,1.45) (4.1,.86) (5.4,.45) (6.6,.22)};
  \foreach \x/\y in {.65/2.2,1.3/2.75,2.0/2.25,3.0/1.45,4.1/.86,5.4/.45,6.6/.22} {
    \fill[pkblue] (\x,\y) circle (.045);
    \draw[draw=pkamber, fill=pklightamber, opacity=.45] (\x,0) -- (\x,\y) -- ++(.001,0) -- (\x,0);
  }
  \draw[pkamber, thick] (.65,2.2) -- (1.3,2.75) -- (2.0,2.25) -- (3.0,1.45) -- (4.1,.86) -- (5.4,.45) -- (6.6,.22);
  \fill[pklightamber, opacity=.55] (.65,0) -- (.65,2.2) -- (1.3,2.75) -- (2.0,2.25) -- (3.0,1.45) -- (4.1,.86) -- (5.4,.45) -- (6.6,.22) -- (6.6,0) -- cycle;
  \draw[pkdash] (6.6,.22) -- (7.85,.05);
  \node[align=center] at (3.8,2.85) {observed AUC:\\trapezoids};
  \node[align=center] at (7.0,.75) {tail:\\\(C_{\pklab{last}}/\lambda_z\)};
\end{tikzpicture}%
}
\captionof{figure}{NCA as curve geometry. The observed area is numerical geometry; the terminal tail is a model assumption about \(\lambda_z\).}
\end{center}
}

\newcommand{\figureMultipleDosing}{
\begin{center}
\pkfitfigure{%
\begin{tikzpicture}[font=\small, x=1cm, y=1cm]
  \draw[->] (0,0) -- (9.1,0) node[right] {time};
  \draw[->] (0,0) -- (0,3.3) node[above] {concentration};
  \foreach \d in {0,1.5,3,4.5,6,7.5} {
    \draw[pkblue, thick, smooth] plot coordinates {(\d,0.15) (\d+.2,2.1) (\d+.65,1.4) (\d+1.2,.8) (\d+1.5,.55)};
    \draw[pkarrred] (\d,3.05) -- (\d,2.35);
  }
  \draw[pkgreen, very thick, smooth] plot coordinates {(0,.15) (.2,2.1) (.65,1.4) (1.5,2.65) (2.15,1.95) (3,2.95) (3.65,2.2) (4.5,3.05) (5.2,2.28) (6,3.1) (6.7,2.33) (7.5,3.12) (8.2,2.35) (9,1.75)};
  \draw[dashed, pkamber, thick] (0,2.35) -- (9,2.35) node[right] {\(\pkCavg\)};
  \draw[<->, thick, draw=pkred] (8.0,2.35) -- (8.0,3.08) node[midway,right] {fluctuation};
  \node at (4.5,-.45) {same interval \(\tau\), accumulation toward a repeating steady-state pattern};
\end{tikzpicture}%
}
\captionof{figure}{Multiple dosing by superposition. Repeated doses accumulate until the profile repeats; average exposure is clearance-controlled, while fluctuation is shaped by half-life and \(\tau\).}
\end{center}
}

\newcommand{\figureWellStirred}{
\begin{center}
\pkfitfigure{%
\begin{tikzpicture}[font=\small, node distance=1cm and 1.1cm]
  \node[pkboxgreen, minimum width=2.1cm, minimum height=1.1cm] (res) {reservoir\\rest of body};
  \node[pkboxred, right=3.0cm of res, minimum width=2.1cm, minimum height=1.1cm] (liv) {well-stirred\\liver};
  \draw[pkarrred] (res.north east) -- node[above] {\(Q_H, C_{\pklab{in}}\)} (liv.north west);
  \draw[pkarr] (liv.south west) -- node[below] {\(Q_H, C_{\pklab{out}}\)} (res.south east);
  \draw[pkarrred] (liv) -- ++(0,-1.35) node[below,align=center] {metabolism\\\(f_u\pkCL_{\pklab{int}}\)};
  \node[pkboxamber, above=1.1cm of res, minimum width=2.2cm] (input) {input\\\(R_0\)};
  \draw[pkarr] (input) -- (res);
  \node[below=1.85cm of liv, align=center] {\(\displaystyle \pkCL_H=\frac{Q_H f_u\pkCL_{\pklab{int}}}{Q_H+f_u\pkCL_{\pklab{int}}}\), \quad \(E_H=(C_{\pklab{in}}-C_{\pklab{out}})/C_{\pklab{in}}\)};
\end{tikzpicture}%
}
\captionof{figure}{Well-stirred hepatic clearance. This redraw keeps the source figure's core idea: blood flow brings drug into a mixed liver space where intrinsic capacity and binding determine extraction.}
\end{center}
}

\newcommand{\figureRenalHandling}{
\begin{center}
\pkfitfigure{%
\begin{tikzpicture}[font=\small, node distance=.9cm and 1.15cm]
  \node[pkboxgreen, minimum width=2.2cm] (blood) {renal blood\\unbound drug};
  \node[pkbox, right=of blood, minimum width=2.2cm] (filtrate) {glomerular\\filtrate};
  \node[pkboxamber, right=of filtrate, minimum width=2.2cm] (tubule) {tubule\\lumen};
  \node[pkboxred, right=of tubule, minimum width=2.0cm] (urine) {urine\\\(A_e\)};
  \draw[pkarr] (blood) -- node[above] {\(f_u\pksym{GFR}\)} (filtrate);
  \draw[pkarr] (filtrate) -- (tubule);
  \draw[pkarr] (tubule) -- (urine);
  \draw[pkarrgreen] (blood.south) to[out=-55,in=-125,looseness=1.05]
    node[pos=.52, above] {secretion} (tubule.south);
  \draw[pkarrred] (tubule.north) to[out=125,in=55,looseness=1.05]
    node[pos=.52, above] {reabsorption} (blood.north);
\end{tikzpicture}%
}
\captionof{figure}{Renal clearance as filtration, secretion, and reabsorption. The drawing deliberately separates tubule lumen from blood so \(\pkCL_R\) can be read mechanistically.}
\end{center}
}

\newcommand{\figureTMDD}{
\begin{center}
\pkfitfigure{%
\begin{tikzpicture}[font=\small, node distance=1.0cm and 1.5cm]
  \node[pkbox] (drug) {free drug\\\(D\)};
  \node[pkboxgreen, right=of drug] (target) {target\\\(R\)};
  \node[pkboxamber, below=1.15cm of target] (complex) {complex\\\(DR\)};
  \draw[pkarr] (drug) -- node[above] {\(k_{\pklab{on}}\)} (target);
  \draw[pkarr] (target) -- (complex);
  \draw[pkarrgreen] (complex.west) to[bend left=18] node[left] {\(k_{\pklab{off}}\)} (drug.south);
  \draw[pkarrred] (complex) -- ++(0,-1.05) node[below,align=center] {internalization\\or target-mediated loss};
  \draw[pkarrred] (drug.south) -- ++(-.95,-.55) node[below,align=center] {linear\\clearance};
  \draw[pkarrgreen] (target.north) -- ++(0,.95) node[above] {target turnover};
\end{tikzpicture}%
}
\captionof{figure}{Target-mediated drug disposition. TMDD makes the pharmacological target part of the disposition model, so PK and PD become coupled rather than sequential.}
\end{center}
}

\newcommand{\figurePopPKHierarchy}{
\begin{center}
\pkfitfigure{%
\begin{tikzpicture}[font=\small, node distance=1.05cm and 1.35cm]
  \node[pkboxamber] (theta) {population\\\(\theta,\Omega,\Sigma\)};
  \node[pkboxgreen, below left=of theta] (cov) {covariates\\\(c_i\)};
  \node[pkbox, below right=of theta] (eta) {random effects\\\(\eta_i\)};
  \node[pkbox, below=2.0cm of theta] (psi) {individual parameters\\\(\psi_i=h(c_i,\theta,\eta_i)\)};
  \node[pkboxred, below left=of psi] (dose) {dose/design\\\(u_i,t_{ij}\)};
  \node[pkboxgreen, below right=of psi] (obs) {observations\\\(y_{ij}\)};
  \draw[pkarr] (theta) -- (eta);
  \draw[pkarr] (theta) -- (psi);
  \draw[pkarr] (cov) -- (psi);
  \draw[pkarr] (eta) -- (psi);
  \draw[pkarr] (dose) -- (obs);
  \draw[pkarr] (psi) -- (obs);
  \node[draw=pkblue, dashed, rounded corners=5pt, fit=(cov)(eta)(psi)(dose)(obs), inner sep=9pt, label={[pkblue]below:subject \(i=1,\ldots,N\)}] {};
\end{tikzpicture}%
}
\captionof{figure}{Population PK hierarchy. NONMEM, SAEM, and Bayesian TDM are different computational cultures around this same latent-individual-parameter graph.}
\end{center}
}

\newcommand{\figurePBPK}{
\begin{center}
\pkfitfigure{%
\begin{tikzpicture}[font=\scriptsize, x=1cm, y=.72cm]
  \node[pkboxred, minimum width=1.7cm] (art) at (0,5) {arterial\\blood};
  \node[pkbox, minimum width=1.7cm] (ven) at (0,0) {venous\\blood};
  \node[pkorgan] (lung) at (3.0,5.0) {lung};
  \node[pkorgan] (brain) at (3.0,4.1) {brain};
  \node[pkorgan] (heart) at (3.0,3.2) {heart};
  \node[pkorgan] (kidney) at (3.0,2.3) {kidney};
  \node[pkorgan] (gut) at (3.0,1.4) {gut};
  \node[pkorgan, minimum width=2.0cm] (liver) at (6.0,1.4) {liver\\metabolism};
  \node[pkorgan] (fat) at (3.0,.5) {fat};
  \node[pkorgan] (muscle) at (3.0,-.4) {muscle};
  \draw[pkarrred] (art) -- (lung);
  \foreach \o/\lab in {brain/\(Q_B\),heart/\(Q_H\),kidney/\(Q_K\),gut/\(Q_G\),fat/\(Q_F\),muscle/\(Q_M\)} {
    \draw[pkarrred] (art.east) -- node[above,sloped] {\lab} (\o.west);
    \draw[pkarr] (\o.west) -- (ven.east);
  }
  \draw[pkarrred] (gut) -- node[above] {portal} (liver);
  \draw[pkarr] (liver) |- (ven.east);
  \draw[pkarr] (ven) -- ++(-1,0) |- (art);
\end{tikzpicture}%
}
\captionof{figure}{A stripped-down PBPK circulation. Whole-body PBPK replaces abstract compartments with blood-flow-connected organs, while still preserving mass balance and observation maps.}
\end{center}
}

\newcommand{\figurePDDelays}{
\begin{center}
\pkfitfigure{%
\begin{tikzpicture}[font=\small, node distance=.9cm and 1.2cm]
  \node[pkbox] (cp) {plasma\\\(C_p(t)\)};
  \node[pkboxgreen, right=of cp] (direct) {direct effect\\\(E(C_p)\)};
  \draw[pkarr] (cp) -- (direct);
  \node[pkboxamber, below=of cp] (ce) {effect site\\\(\dot C_e=k_{e0}(C_p-C_e)\)};
  \node[pkboxgreen, right=of ce] (effect) {delayed effect\\\(E(C_e)\)};
  \draw[pkarr] (cp) -- (ce);
  \draw[pkarr] (ce) -- (effect);
  \node[pkboxred, below=of ce, minimum width=3.1cm] (turn) {turnover response\\\(\dot R=k_{\pklab{in}}S(C)-k_{\pklab{out}}R\)};
  \draw[pkarr] (cp) -- (turn);
  \node[align=left, right=1.1cm of turn] {use when the response state\\has its own production/loss\\time scale};
\end{tikzpicture}%
}
\captionof{figure}{Three PD delay grammars. Direct effect, effect compartment, and turnover models answer different reasons why response and plasma concentration may disconnect.}
\end{center}
}

\newcommand{\figureStructureAbsorptionMap}{
\begin{center}
\pkfitfigure{%
\begin{tikzpicture}[font=\small, node distance=.85cm and 1.0cm]
  \node[pkboxamber, minimum width=2.3cm] (mol) {molecular\\structure};
  \node[pkbox, right=of mol, minimum width=2.35cm] (desc) {descriptors\\\(pK_a,\log P\)\\PSA, MW};
  \node[pkboxgreen, right=of desc, minimum width=2.5cm] (hyp) {ADME\\hypotheses};
  \node[pkboxred, right=of hyp, minimum width=2.45cm] (model) {PK model\\components};
  \node[pkboxgreen, right=of model, minimum width=2.3cm] (data) {curve and\\decision};
  \draw[pkarr] (mol) -- (desc);
  \draw[pkarr] (desc) -- (hyp);
  \draw[pkarr] (hyp) -- (model);
  \draw[pkarr] (model) -- (data);
  \node[pkboxamber, below=of hyp, minimum width=6.6cm] (examples) {solubility / dissolution \quad permeability / efflux \quad first pass \quad tissue partitioning};
  \draw[pkdash] (examples) -- (hyp);
\end{tikzpicture}%
}
\captionof{figure}{Structure-PK as hypothesis generation. Molecular descriptors do not replace clinical data; they propose where to look in absorption, distribution, metabolism, excretion, and tissue exposure.}
\end{center}
}

\newcommand{\figureReversibleEhc}{
\begin{center}
\pkfitfigure{%
\begin{tikzpicture}[font=\small, node distance=.85cm and 1.05cm]
  \node[pkboxgreen, minimum width=2.05cm] (central) {central\\plasma};
  \node[pkboxred, right=of central, minimum width=2.05cm] (liver) {liver};
  \node[pkboxamber, right=of liver, minimum width=2.15cm] (bilegut) {bile / gut\\loop};
  \node[pkboxgreen, below=of central, minimum width=2.05cm] (met) {metabolite};
  \draw[pkarr] (central) -- node[above] {\(Q_H\)} (liver);
  \draw[pkarr] (liver) -- node[above] {biliary} (bilegut);
  \draw[pkarrgreen] (bilegut.south) to[out=-90,in=-90,looseness=1.15] node[pos=.5,below] {reabsorption} (central.south);
  \draw[pkarrred] (bilegut.east) -- ++(.9,0) node[right] {fecal loss};
  \draw[pkarr] (central) -- (met);
  \draw[pkarrgreen] (met.west) to[out=180,in=-90,looseness=1.05] node[pos=.55,left] {reversible} (central.west);
\end{tikzpicture}%
}
\captionof{figure}{Reversible systems and enterohepatic cycling. A secondary peak may be a real loop in the system, not merely noisy concentration data.}
\end{center}
}

\newcommand{\figureAdvancedPdDelayLadder}{
\begin{center}
\pkfitfigure{%
\begin{tikzpicture}[font=\scriptsize, node distance=.52cm and .55cm]
  \node[pkbox, minimum width=1.55cm] (direct) {direct\\effect};
  \node[pkboxgreen, right=of direct, minimum width=1.55cm] (bio) {biophase\\delay};
  \node[pkboxamber, right=of bio, minimum width=1.7cm] (turnover) {turnover\\IRM};
  \node[pkboxred, right=of turnover, minimum width=1.75cm] (trans) {transduction\\chain};
  \node[pkboxgreen, right=of trans, minimum width=1.75cm] (life) {cell\\lifespan};
  \node[pkboxamber, right=of life, minimum width=1.75cm] (disease) {disease\\progression};
  \draw[pkarr] (direct) -- (bio);
  \draw[pkarr] (bio) -- (turnover);
  \draw[pkarr] (turnover) -- (trans);
  \draw[pkarr] (trans) -- (life);
  \draw[pkarr] (life) -- (disease);
  \node[below=.65cm of turnover, align=center] {add delay only when a mechanism, diagnostic pattern, or decision needs it};
\end{tikzpicture}%
}
\captionof{figure}{The PHC 609 delay ladder. Advanced PD is not one model class; it is a disciplined escalation from immediate effect to biophase, turnover, transduction, cell lifespan, and disease dynamics.}
\end{center}
}

\newcommand{\figureQSPWorkflow}{
\begin{center}
\pkfitfigure{%
\begin{tikzpicture}[font=\small, node distance=.75cm and .9cm]
  \node[pkboxamber, minimum width=2.2cm] (q) {question\\of interest};
  \node[pkbox, right=of q, minimum width=2.2cm] (rxn) {reaction\\network};
  \node[pkboxgreen, right=of rxn, minimum width=2.2cm] (ode) {\(\dot x=Sv(x)\)\\or stochastic};
  \node[pkbox, right=of ode, minimum width=2.2cm] (sim) {SBML /\\Tellurium};
  \node[pkboxred, right=of sim, minimum width=2.2cm] (check) {sensitivity\\and checks};
  \draw[pkarr] (q) -- (rxn);
  \draw[pkarr] (rxn) -- (ode);
  \draw[pkarr] (ode) -- (sim);
  \draw[pkarr] (sim) -- (check);
  \draw[pkdash] (check.south) to[out=-90,in=-90,looseness=1.15] node[pos=.5,below] {revise mechanism or decision} (q.south);
\end{tikzpicture}%
}
\captionof{figure}{QSP implementation loop. A pathway diagram becomes useful when it is turned into equations, simulated reproducibly, stress-tested, and tied back to a decision.}
\end{center}
}
